% Fragmento público generado desde propietarios íntegros.
% Residencia: 52.3.2.
% Fuente: ../incoming/radion_monografia_20260827/manuscrito/sections/04_cincuenta_critica.tex:106-171
\subsubsection{Separación tipada de los 2 objetos de valor \(1/2\)}

APP contiene además un modo singular con
\[
s=\frac12,\qquad s^2=\frac14,\qquad
1-s^2=\frac34=\frac{90}{120},\qquad
s\sqrt{1-s^2}=\frac{\sqrt3}{4}.
\]
Este \(s\) nace de los ángulos principales de las hojas suma y producto. La
costura crítica \(\Re\rho=1/2\) nace de una involución funcional. Que ambas
coordenadas valgan \(1/2\) es una coincidencia escalar con posible contenido,
pero no autoriza a identificarlas sin un cuadrado tipado.

La forma correcta de conservar el hallazgo es un diagrama de igualadores:
\[
\begin{tikzpicture}[baseline=(current bounding box.center),node distance=14mm and 23mm,
 every node/.style={font=\small}]
\node[draw=RadBlue,rounded corners,fill=RadMist] (crit) {costura crítica \(\Re s=1/2\)};
\node[draw=RadTeal,rounded corners,fill=RadCream,right=of crit] (face) {carta \(j_{100}\)};
\node[draw=RadCopper,rounded corners,fill=RadCream,right=of face] (theta) {fase \(\theta_2=50^\circ\)};
\node[draw=RadRose,rounded corners,fill=white,below=of face] (mode) {modo APP \(s=1/2\), \(1-s^2=90/120\)};
\draw[-{Latex},thick,RadBlue] (crit)--node[above]{\(j_{100}\)}(face);
\draw[-{Latex},thick,RadTeal] (face)--node[above]{igualador}(theta);
\draw[-{Latex},dashed,RadRose] (mode)--node[right,align=left]{misma cifra;\\dominio distinto}(face);
\end{tikzpicture}
\]
El diagrama no finge la flecha ausente. Muestra qué está demostrado, qué se
ha formalizado y qué igualdad sigue necesitando un transporte adicional si
se pretende una identificación operatoria global.

\subsubsection{Compatibilidad entre círculo espectral y memoria helicoidal}

El lector Cayley--Pontryagin transporta una fibra espectral mediante
\[
\widetilde U_{\Gamma_9^n}J\psi_\gamma
=\tau_\gamma^nJ\psi_\gamma.
\]
La publicación circular es \(\tau_\gamma^n\in S^1\). El lift
\[
(\tau_\gamma^n,n)\in S^1\times\Z
\]
conserva el número de retornos. En la hoja contractiva aparece
\[
\widetilde M_9^nJ\psi_\gamma
=\left(\frac59\right)^n\tau_\gamma^nJ\psi_\gamma.
\]
Tenemos, por tanto, tres imágenes tipadas:
\[
\text{círculo de fase},\qquad
\text{hélice de memoria},\qquad
\text{espiral contractiva}.
\]

Un cero no es una vuelta, y un nueve no es un cero particular de la función
zeta. El cero fija un carácter \(\tau_\gamma\) que persiste a través de las
vueltas; el retorno nonádico incrementa memoria. El nueve funciona como cero
residual en \(\Z/9\Z\), una operación aritmética distinta de la enumeración
de ceros espectrales.

\begin{narrativebox}
La recta crítica se vuelve círculo cuando olvidamos la altura lineal y
conservamos el carácter unitario. Se vuelve hélice cuando volvemos a insertar
la memoria de cuántas veces el carácter ha sido transportado. La profundidad
no es una tercera coordenada decorativa: es la información que la proyección
circular no puede retener.
\end{narrativebox}
